\section{Reference Reuse and Follow-up Design}
The original annotation plan was amended before the follow-up evaluations to reuse same-version published NL4Opt references. The original \HeldoutN{} IDs and expanded \ReferenceN{} IDs are preserved, with nested cohorts, source hashes, record locations, exact parsing, feasibility, and bounded-optimum checks. Known corrections remain in the separate sensitivity branch. Blank human packages remain archived; no human reading-status labels were collected. Krippendorff alpha is inapplicable to a single-source reference import. One auxiliary model consultation reviewed the protocol without producing per-case annotation labels.

The follow-ups use DeepSeek flash at temperature zero, with continuous nonnegative variables. Each paid stage records its sample size, arms, switching rules, queue hash, and cost estimate before requests. Failures and invalid responses are retained without retries. Selected cohorts include public training and development problems; their scores describe these selected cohorts rather than an independent benchmark test split.

\subsection{Repeated Sampling (T7)}
Six identical-prompt draws per problem are clustered by feasible-region and optimal-face equivalence; the modal class supplies the selected answer. The cohort is deterministically selected from eligible representatives before new outcomes. Temperature zero still permits different observed responses. Table~\ref{tab:consistency} reports modal success and conditional error, retaining incomplete cases separately. Mean modal share is \ModalShare{}\%, and mean within-problem exact-text repetition is \TextRepeat{}\%; problem bootstrap intervals are in the numerical source. Agreement remains compatible with a shared incorrect formulation.

\subsection{Feedback Granularity (T8)}
Screening uses the original two distinct A/B prompts; T7's identical-prompt draws do not define these strata. Same-version train/dev data supplement local cases. The eligible pool contains \ScreenN{} problems. Predeclared hash ordering selects \DivergentN{} divergent pairs and \WrongAgreeN{} agreeing-but-wrong pairs. Source counts are \LadderTrain{} train, \LadderDev{} development, and \LadderTest{} test problems.

Each case starts from A, keeps B fixed, and receives verdict-only, point, or point-plus-anchor feedback for at most three rounds. Anchors are verbatim source spans retrieved heuristically using official constraint annotations; missing anchors remain missing. This assistance is part of the experimental condition. The official formula and correctness label are not supplied in prompts. A case-arm stops after its first reference-equivalent output or an invalid response. This stopping rule uses reference access and defines an evaluation procedure.

Success@3 means any successful post-feedback output, with every assigned problem retained in its arm denominator. It includes initially correct candidates. The planning MDE is approximately \PlanningMDE{} percentage points at $n=\LadderN{}$; it is not a treatment-effect confidence bound. Table~\ref{tab:ladder} includes both strata and exact marginal intervals. The largest observed arm difference is \ObservedGap{} percentage points, below that planning threshold. Of \InitiallyWrong{} initially wrong A candidates, the arms repair \VerdictRepair{}, \PointRepair{}, and \AnchorRepair{}; of \InitiallyCorrect{} initially correct candidates, \VerdictLost{}, \PointLost{}, and \AnchorLost{} have no successful post-feedback output. These are descriptive transitions.

Table~\ref{tab:rounds} records compliance (membership of the preceding witness in the new feasible set) and damage (loss of a reference row previously implied by the nonempty candidate region). Witness acceptance and correct membership are separate: rejecting a reference-infeasible point can be correct. The verdict arm's witness is measured privately, not exposed. Undefined metrics retain null values and their reasons. First-round anchors reach \AnchorExposed{} cases. The later rounds add zero successes; all valid third-round outputs equal their preceding candidates in the contract.

\subsection{Metamorphic Alarms (T11)}
The amended criterion evaluates alarms relative to published references, with no human false-positive estimate. Pre-outcome hash selection takes \MetaN{} distinct problems supporting all three transformations from \SupportedN{} eligible triples. Units count groups of ten original units; annotated variable mentions become item A/B; explicit percentage/multiplier literals are rewritten using exact rational ratios. Unsupported expressions, including ``times more,'' are excluded by rule. Unit scaling maps $x=10z$; the other transformations preserve coordinates. Existing valid A outputs are reused as baselines. Source counts are \MetaTrain{} train, \MetaDev{} development, and \MetaTest{} test problems.

An alarm flags a transformed answer when it differs from the mapped baseline under the contract. A reference-conditional false alarm is an alarm whose transformed answer matches the mapped official model. Its denominator is all valid alarms: this is a false-discovery fraction, not the population false-positive rate. An initially wrong baseline followed by a correct transformed answer can cause it. Text meaning preservation is rule-based and was not independently human-validated.

Table~\ref{tab:metamorphic} retains valid-pair denominators and invalid outputs. Pooled reference-conditional false alarms are \FalseK{}/\FlagN{} (\FalseRate{}\%), exceeding the preregistered \DiscardThreshold{}\% threshold; the module is discarded. At least one violation occurs on \AnyViolationK{}/\CompleteMetaN{} complete problems. Transformations of the same problem are correlated, so pooled alarms do not receive an independent-binomial interval.

\begin{table*}[t]
\centering\small
\begin{tabular}{lrrl}
\hline
T7 outcome & $k/n$ & Rate (\%) & Exact 95\% CI (\%) \\
\hline
All: modal success & 76/100 & 76.00 & 66.43--83.98 \\
Fully agreeing: error & 13/83 & 15.66 & 8.61--25.29 \\
Divergent: error & 7/12 & 58.33 & 27.67--84.83 \\
Incomplete: error & 4/5 & 80.00 & 28.36--99.49 \\
\hline
\end{tabular}
\caption{Repeated sampling. All numerical cells derive from ../results.json, consistency.majority\_success and consistency.conditional\_errors. Incomplete cases remain in the overall denominator.}
\label{tab:consistency}
\end{table*}
\begin{table*}[t]
\centering\small
\begin{tabular}{llrrl}
\hline
Stratum & Arm & Success/$n$ & Rate (\%) & Exact 95\% CI (\%) \\
\hline
All & verdict & 19/80 & 23.75 & 14.95--34.58 \\
All & point & 20/80 & 25.00 & 15.99--35.94 \\
All & point+anchor & 23/80 & 28.75 & 19.18--39.95 \\
Divergent & verdict & 19/60 & 31.67 & 20.26--44.96 \\
Divergent & point & 20/60 & 33.33 & 21.69--46.69 \\
Divergent & point+anchor & 23/60 & 38.33 & 26.07--51.79 \\
Agreeing but wrong & verdict & 0/20 & 0.00 & 0.00--16.84 \\
Agreeing but wrong & point & 0/20 & 0.00 & 0.00--16.84 \\
Agreeing but wrong & point+anchor & 0/20 & 0.00 & 0.00--16.84 \\
\hline
\end{tabular}
\caption{T8 success within at most three rounds. Every numerical cell derives from ../results.json, ladder.main\_table. Marginal intervals use distinct problems within each arm; they are not paired arm-difference intervals.}
\label{tab:ladder}
\end{table*}
\begin{table*}[t]
\centering\small
\begin{tabular}{lrrll}
\hline
Arm & Round & Requests & Compliance $k/n$ (CI\%) & Damage $k/n$ (CI\%) \\
\hline
verdict & 1 & 80 & 34/60 (43.2--69.4) & 7/77 (3.7--17.8) \\
verdict & 2 & 61 & 10/21 (25.7--70.2) & 0/59 (0.0--6.1) \\
verdict & 3 & 61 & 10/20 (27.2--72.8) & 0/59 (0.0--6.1) \\
point & 1 & 80 & 32/60 (40.0--66.3) & 7/76 (3.8--18.1) \\
point & 2 & 59 & 10/18 (30.8--78.5) & 1/57 (0.0--9.4) \\
point & 3 & 59 & 8/16 (24.7--75.3) & 0/57 (0.0--6.3) \\
point+anchor & 1 & 80 & 35/59 (45.7--71.9) & 8/75 (4.7--19.9) \\
point+anchor & 2 & 55 & 8/14 (28.9--82.3) & 0/53 (0.0--6.7) \\
point+anchor & 3 & 55 & 8/14 (28.9--82.3) & 0/53 (0.0--6.7) \\
\hline
\end{tabular}
\caption{T8 round measurements, with exact marginal intervals. Cells derive from ../results.json, ladder.main\_table.all.*.rounds. Denominators exclude undefined metrics; stopped trajectories remain in the success denominator.}
\label{tab:rounds}
\end{table*}
\begin{table*}[t]
\centering\small
\begin{tabular}{lrrlrrr}
\hline
Transform & Violations/$n$ & Rate (\%) & Exact 95\% CI (\%) & False/alarms & Fraction (\%) & Invalid \\
\hline
Unit scaling & 139/294 & 47.28 & 41.45--53.16 & 18/139 & 12.95 & 6 \\
Variable rename & 64/297 & 21.55 & 17.01--26.67 & 20/64 & 31.25 & 3 \\
Percentage/ratio & 84/298 & 28.19 & 23.15--33.66 & 21/84 & 25.00 & 2 \\
\hline
\end{tabular}
\caption{T11 reference-conditional results. Cells derive from ../results.json, metamorphic.by\_transformation. False alarms count reference-correct transformed answers among valid flags. Per-transform false-alarm intervals are also retained in the numerical source.}
\label{tab:metamorphic}
\end{table*}


\section{Reproducibility}
The facts sheet maps table cells and narrative macros to source fields. The public manuscript directory includes selected numerical outcomes, source hashes, and typesetting inputs. Raw prompts, responses, credentials, and source snapshots remain local under the repository's publication policy. Public checks reproduce derived aggregations; full geometric replay requires the archived raw inputs.
